\section{Background and Related Work}

\subsection{Multimodal Vision-Language Models}
Vision-language models connect a vision encoder to a language model through a lightweight projector, enabling end-to-end visual understanding and instruction following.
% 中文：多模态视觉-语言模型。视觉-语言模型通过轻量投影层把视觉编码器连接到语言模型，实现端到端视觉理解与指令遵循。
LLaVA~\cite{liu2023llava} established visual instruction tuning---aligning a vision encoder to a frozen LLM via a linear projection followed by instruction fine-tuning---as a scalable recipe for open multimodal assistants.
% 中文：LLaVA \cite{liu2023llava} 把视觉指令微调——通过线性投影将视觉编码器对齐到冻结的 LLM，再进行指令微调——确立为开放多模态助手的可扩展配方。
LLaVA-OneVision~\cite{li2024llavaonevision} consolidates this line into a single open model spanning single-image, multi-image, and video tasks through a staged curriculum: stage-1 projector pre-training, stage-1.5 mid-training, and stage-2 instruction tuning.
% 中文：LLaVA-OneVision \cite{li2024llavaonevision} 将这一路线整合为单一开放模型，通过分阶段课程覆盖单图、多图与视频任务：stage-1 投影层预训练、stage-1.5 中期训练、stage-2 指令微调。
BaiZe follows this curriculum in stage~(iv) (Section~\ref{sec:mllm}), substituting a from-scratch vision encoder (stage~iii) for the frozen off-the-shelf encoder.
% 中文：BaiZe 在第 (iv) 阶段沿用这一课程（第 \ref{sec:mllm} 节），但把冻结的现成编码器替换为从零训练的视觉编码器（第 iii 阶段）。

\subsection{LLMs for Electronic Design Automation}
A growing line of work applies LLM agents to chip design tasks such as register-transfer-level generation, verification, and bug triage~\cite{liu2023chipnemo,blocklove2024chip}.
% 中文：面向电子设计自动化的 LLM。越来越多的工作把 LLM 智能体应用于芯片设计任务，如寄存器传输级生成、验证与缺陷分诊 \cite{liu2023chipnemo,blocklove2024chip}。
These applications motivate a cheap, always-on, agentic multimodal core that can run per-team or per-engineer; our work supplies the efficient multimodal backbone such systems require.
% 中文：这些应用呼唤一种廉价、常驻、可按团队或按工程师部署的智能体多模态内核；我们的工作正是提供此类系统所需的、高效的多模态骨干网络。
